\chapter{Bronchoscopy}

\section*{Definition and Overview}
Bronchoscopy is a procedure in which a thin, flexible tube fitted with a camera and light at its tip is passed through the mouth or nose and down into the airways of the lungs. It enables a physician to examine directly the trachea and bronchi, to collect samples of tissue, fluid, or cells, and occasionally to deliver treatment. Most bronchoscopies are performed as day procedures under light sedation and local anaesthesia to the throat. The procedure is commonly used to investigate a persistent cough, coughing up blood, hoarseness, or abnormal findings on a chest scan, and it is an essential tool in both the diagnosis and management of respiratory disease.

\section*{Epidemiology}
Bronchoscopy is one of the most frequently performed diagnostic procedures in respiratory medicine, with hundreds of thousands of procedures carried out annually in large countries. It is used across all ages but is most often performed in adults, particularly those over 50, in whom lung cancer, chronic lung disease, and persistent respiratory symptoms are more common. The number of procedures has grown with the increasing use of CT scanning, which detects suspicious lung nodules and masses that require tissue diagnosis. Bronchoscopy is performed in both inpatient and outpatient settings and is increasingly guided by real-time imaging techniques that improve sampling accuracy.

\section*{Aetiology and Pathophysiology}
Bronchoscopy is not a disease and has no underlying causes of its own; it is a diagnostic and therapeutic tool for conditions affecting the airways and lungs. The diseases it is used to investigate span infection, inflammation, malignancy, and mechanical obstruction.

\begin{itemize}
    \item \textbf{Diagnostic reasons:} investigating a chronic or unexplained cough, haemoptysis, hoarseness, unexplained breathlessness, or an abnormal chest X-ray or CT scan
    \item \textbf{Sampling:} obtaining bronchoalveolar lavage (fluid washes), bronchial brushings, or endobronchial and transbronchial biopsy specimens for microscopy, culture, and pathology
    \item \textbf{Therapeutic reasons:} removing a foreign body, clearing retained secretions, controlling bleeding, or placing a stent to keep a narrowed airway open
    \item \textbf{Staging:} assessing the extent of lung cancer, including sampling of lymph nodes via endobronchial ultrasound (EBUS) where available
\end{itemize}

\section*{Clinical Features}
The reasons a person is referred for bronchoscopy are the clinical features of the underlying respiratory condition.

\begin{itemize}
    \item A persistent or unexplained cough lasting more than a few weeks
    \item Coughing up blood (haemoptysis), even in small amounts
    \item Unexplained breathlessness or a change in breathing pattern
    \item Hoarseness or a change in the voice
    \item A lung nodule, mass, or area of collapse identified on chest imaging
    \item Recurrent pneumonia or an infection that fails to clear with treatment
    \item Constitutional symptoms such as weight loss or fever when malignancy is suspected
\end{itemize}

\section*{Differential Diagnosis}
The symptoms that prompt bronchoscopy arise from many possible conditions, and the procedure is often needed specifically to distinguish between them.

\begin{itemize}
    \item \textbf{Lung cancer:} primary bronchogenic carcinoma or metastatic disease
    \item \textbf{Infection:} tuberculosis, bacterial or viral pneumonia, or fungal lung disease
    \item \textbf{Inflammatory lung disease:} sarcoidosis, interstitial lung disease, or allergic bronchopulmonary aspergillosis
    \item \textbf{Foreign body:} an object lodged within the airway
    \item \textbf{Chronic airways disease:} chronic bronchitis, bronchiectasis, or asthma
    \item \textbf{Upper-airway bleeding:} haemorrhage from the nose or throat mistaken for a pulmonary source
\end{itemize}

\section*{When to Seek Medical Care}
You may be referred for bronchoscopy because of an abnormal scan or persistent symptoms; report any worrying symptom to your doctor promptly. Before the procedure, tell the team about all medications you take, particularly blood thinners, and any allergies or significant health conditions.

\textbf{Call 000 (or local emergency services) for:}
\begin{itemize}
    \item Severe difficulty breathing after the procedure
    \item Coughing up a large amount of blood
    \item Chest pain or a rapid heartbeat
    \item High fever or chills within a day or two of the procedure
    \item Difficulty swallowing, or swelling of the throat
\end{itemize}

\section*{Diagnosis and Investigations}
\begin{itemize}
    \item \textbf{Pre-procedure history and examination:} review of current symptoms, medications, and general fitness for sedation
    \item \textbf{Chest imaging:} a chest X-ray or CT scan performed beforehand to guide the areas examined during the procedure
    \item \textbf{Blood tests:} a full blood count and clotting studies before any planned biopsy
    \item \textbf{The procedure itself:} the bronchoscope is passed while the patient is relaxed and sedated, and the airways are inspected directly
    \item \textbf{Specimen analysis:} washings, brushings, and biopsies are sent to the laboratory for cytology, microbiology, and histopathology
    \item \textbf{EBUS sampling:} where indicated, ultrasound-guided needle sampling of mediastinal lymph nodes to stage lung cancer
\end{itemize}

\section*{Management}
Bronchoscopy is most often a diagnostic step, but it can also treat certain conditions.

\begin{itemize}
    \item \textbf{Preparation:} fast for several hours beforehand, arrange transport home, and follow the sedation instructions provided
    \item \textbf{Recovery:} rest until the effects of sedation wear off; a mild sore throat and a little blood-streaked sputum are common
    \item \textbf{Sample-guided therapy:} antibiotics, antifungals, or anticancer treatment selected according to the laboratory results
    \item \textbf{Removal of a foreign body:} small objects can usually be grasped and removed during the procedure
    \item \textbf{Airway interventions:} clearance of thick secretions, laser or debulking treatment of lesions, or stenting of a narrowed airway
    \item \textbf{Follow-up:} discussion of results with a specialist and planning of any further tests or treatment
\end{itemize}

\section*{Complications}
Bronchoscopy is generally safe, and serious complications are uncommon.

\begin{itemize}
    \item Sore throat, hoarseness, or minor bleeding after biopsy
    \item Temporary lowering of blood oxygen or an irregular heart rhythm during the procedure
    \item Pneumothorax (collapsed lung) if a biopsy punctures the lung surface
    \item Heavier-than-expected bleeding from a biopsy site
    \item An adverse reaction to the sedative or local anaesthetic
    \item Rarely, infection introduced by the procedure
\end{itemize}

\section*{Prognosis and Outcomes}
Most people tolerate bronchoscopy very well and return to normal activities within a day. The outlook depends entirely on the underlying condition the procedure reveals: a benign finding such as infection usually responds to treatment, whereas the prognosis of lung cancer depends on its type and stage. When bronchoscopy is performed to remove a foreign body or clear an obstructed airway, symptoms often improve rapidly. Results should be discussed with the treating specialist, who will explain their significance and recommend next steps.

\section*{Prevention and Self-Care}
Because bronchoscopy is a procedure rather than an illness, prevention centres on the conditions that make it necessary.

\begin{itemize}
    \item Do not smoke, and avoid tobacco smoke and inhaled irritants that damage the lungs
    \item See a doctor early for a persistent cough, coughing up blood, or unexplained breathlessness
    \item Follow the fasting and medication instructions given before the procedure
    \item Arrange for someone to drive you home and stay with you after sedation
    \item Protect lung health with influenza and pneumococcal vaccination if recommended
    \item Report any unusual symptoms after the procedure without delay
\end{itemize}

\section*{Sources and Further Reading}
The following organisations provide reliable information on bronchoscopy for patients and healthcare students.

\begin{itemize}
    \item NHS. \textit{Bronchoscopy: what happens during the procedure and recovery.} https://www.nhs.uk/conditions/bronchoscopy/
    \item Mayo Clinic. \textit{Bronchoscopy: purpose, risks, and what to expect.} https://www.mayoclinic.org/
    \item MedlinePlus. \textit{Bronchoscopy: preparation, procedure, and results.} https://medlineplus.gov/bronchoscopy.html
    \item MSD Manual. \textit{Bronchoscopy: professional overview and indications.} https://www.msdmanuals.com/
    \item American Thoracic Society. \textit{Patient information on bronchoscopy and related procedures.} https://www.thoracic.org/
    \item Thoracic Society of Australia and New Zealand. \textit{Respiratory procedure information for patients and clinicians.} https://www.thoracic.org.au/
\end{itemize}
